% SPDX-License-Identifier: Apache-2.0
\section{Method}\label{sec:method}

\paragraph{The rule.} With a classical scorer $f$ and thresholds $\taulo < \tauhi$ fixed in
advance,
\[
D(x) = \begin{cases}
\textsc{decline} & f(x) \ge \tauhi\\
\textsc{step-up}, \text{ then } g(x) \ge \lambda \Rightarrow \textsc{decline} & \taulo < f(x) < \tauhi\\
\textsc{approve} & f(x) \le \taulo
\end{cases}
\]
The band is $B = \{x : \taulo < f(x) < \tauhi\}$ and $g$ is the in-band re-scorer, which is
where a quantum model would enter if one were usable.

\paragraph{The estimand is band-conditional, deliberately.} We certify
\[
R(\lambda) \;=\; \mathbb{P}\bigl(D(X) = \textsc{decline} \,\bigm|\, Y = 0,\; X \in B\bigr),
\]
not the marginal false-decline rate. The marginal rate is dominated by $\tauhi$ and would be
almost unchanged if $g$ were replaced by a coin flip. A certificate that does not bind the
component it is said to license is not evidence about that component --- so the harder,
smaller-sample quantity is the one we bound, and the cost of that choice is reported in
Section~\ref{sec:results} rather than hidden.

\paragraph{Four blocks, and two distinct breaks of exchangeability.} $\Dtrain$ (60\,\%) fits
$f$; $\Dband$ (10\,\%) is the \emph{only} data that may set the band edges, $g$'s
hyperparameters or any screen threshold; $\Dcal$ (10\,\%) certifies the composite rule;
$\Dtest$ (20\,\%) is read once. Blocks are contiguous and snapped to day boundaries.

Two breaks must not be conflated. The \textbf{temporal} break is that $\Dcal$ precedes
$\Dtest$. The \textbf{selective} break is that conditioning on band membership is
conditioning on a data-dependent event --- if the edges were estimated on the same data used
to certify. Freezing the band on $\Dband$ removes the second by construction: membership in
$B$ becomes a fixed measurable predicate, and filtering a sequence by a fixed predicate
preserves exchangeability of what remains. Only the temporal break survives, and we measure
it.

\paragraph{Risk control.} Learn-then-Test controls family-wise error over a finite grid fixed
in advance, using Hoeffding--Bentkus $p$-values with Holm. Two risks are controlled jointly:
$R(\lambda)$ at level $\alpha$, and the false-negative rate at $\alpha_{\mathrm{FN}}$.

The bound applies only to a \emph{mean of per-observation losses}. Our first formulation used
$\max(0,\ \text{floor} - \text{recall})/\text{floor}$, a nonlinear transform of a mean; the
quantities it produced were not valid $p$-values and the certificate resting on them was not
a certificate. It is now a genuine $0/1$ loss averaged over frauds, which is equivalent to a
recall floor at $1 - \alpha_{\mathrm{FN}}$.

\paragraph{Coverage is judged against the exact law, not against $\alpha$.} For a sound
split-conformal system the test-error count follows
$E \sim \mathrm{BetaBinomial}(m,\ n{+}1{-}k,\ k)$ with $k = \lceil (n{+}1)(1-\alpha) \rceil$.
Checking $\hat r \le \alpha$ is a one-sided test against the wrong null: on a correctly
calibrated system it passes only about half the time. Every verdict below is a two-sided
Beta-Binomial interval.

\paragraph{Intervals resample entities, not rows.} The IEEE-CIS label propagates a reported
chargeback to later transactions sharing an account, email or billing address, so rows are
clustered: \ClaimCardOverlap{} of test-block \texttt{card1} values also appear in training.
Every interval is a card-level block bootstrap.
